\section{Horizon complexity and the price of accuracy}\label{sec:complexity86}
The inherited recursive query experiment exposes descendant growth. The own-future constructive backend has a different dependency graph: once $f_{t+1}$ is stored, a query evaluates that representation and does not solve the remaining Bellman problem again. This distinction changes the horizon account without removing any of the recursive benchmarks.

For each date let $S_t$ be its state-context cover size, $J_{ti}$ its feasible action-net size at node $i$, $M_t$ its certified integration work in primitive representatives, and $E_{t+1}$ the cost of evaluating its stored successor. Construction uses
\begin{equation}\label{eq:dagwork86}
 Q_t=M_t\sum_{i\le S_t}J_{ti},\qquad
 W\le\sum_{t<T}\{Q_t(C_t^{\rm prim}+E_{t+1})+C_t^{\rm compile}\}
       +C_T^{\rm terminal}+C^{\rm output}.
\end{equation}
The formula counts a directed acyclic sequence of stored continuations. A literal descendant solver instead multiplies successor solves by its action and shock branching factors. Neither expression is a lower bound for every possible algorithm. The parameters in \eqref{eq:dagwork86} may themselves grow with horizon if the primitive regularity deteriorates.

\begin{proposition}[Optimal dated error allocation]\label{prop:allocation86}
Suppose a certified constructive implementation has a resource upper account $a_t\eta_t^{-p}$ for attaining total date-$t$ residual allowance $\eta_t>0$, with common $p>0$. Each $a_t$ includes action search, integration, feasibility, compilation and the cost of evaluating the stored successor; it cannot hide another recursive Bellman solve. Put $w_t=B_t>0$ and
\[
 S=\sum_t a_t^{1/(p+1)}w_t^{p/(p+1)}.
\]
Among allocations satisfying $\sum_tw_t\eta_t\le\varepsilon$, the minimum of this resource upper model and its minimizer are
\begin{equation}\label{eq:allocation86}
 W_{\rm model}^*=S^{p+1}\varepsilon^{-p},\qquad
 \eta_t^*=\frac{\varepsilon}{S}(a_t/w_t)^{1/(p+1)}.
\end{equation}
This optimizes the declared upper model, not the unknown fastest algorithm. Rounding each allowance down to a dyadic number within a factor two preserves the error constraint and increases that model by at most $2^p$.

If $a_t\le A$ and $w_t=\beta^t$ for $0<\beta<1$, then
\[
 W_{\rm model}^*\le
 A\varepsilon^{-p}(1-\beta^{p/(p+1)})^{-(p+1)}.
\]
With a maximum permissible allowance $\eta_{\max}$ and positive per-date setup costs, clipping the displayed allocation adds at most
$\sum_t a_t\eta_{\max}^{-p}$ to this upper bound, plus the explicit setup account. For $w_t=1$, the un-clipped bound is $A(T+1)^{p+1}\varepsilon^{-p}$. Thus uniform regularity and reusable stored continuations give polynomial, rather than descendant-tree, horizon dependence at fixed dimensions. They do not give dimension-free computation.
\end{proposition}

The power model follows, for example, when effective state, action and innovation covers and their certified errors have polynomial work bounds, successor evaluation has the separately proved compiled bound, and all moduli are uniformly controlled. The exponent must include all these factors. Tensor compilation in Theorem~\ref{thm:compile49} has a $2^{p_t}$ query factor and cell-location work; a dense neural cone has an $S_{t+1}$ factor. Those are different implementation accounts and must not be interchanged.

\subsection{Operation vectors, arithmetic and reliability}\label{sec:ledger86}
We retain the resource vector
\[
 (N_{\rm primitive},N_{\rm descendant},N_{\rm interval},N_{\rm rational},
 N_{\rm network},B_{\rm read},B_{\rm write},M_{\rm peak}).
\]
Counts are reported before scalar prices are applied. Rational operations additionally record operand bit lengths; an arithmetic-operation count is not a bit-complexity bound. Multiplication and division costs at those lengths enter a bit-cost account. Loading, validation, serialization and shutdown belong to complete process time. Peak memory is a capacity requirement, not an additive flow cost unless its duration or rental convention is specified.

A fitting warning, invalid proposal, zero-progress probe, successful screen, protected fallback and terminal service status are different events. In particular, zero new fits cannot estimate fitting reliability, and a successful fallback cannot be counted as a successful prediction. The retained 336-service recursive catalogue and the nine-controller attained-cost experiment have different units of observation. New mathematical fixtures below establish the operation identities and allocation formula; they are not relabeled as a re-execution of the inherited eight-million-query vector tree or as a hardware calibration. The complete proof of Proposition~\ref{prop:allocation86} is in Supplement~\ref{supp:allocation86}.
